\begin{abstract}
Somatic structural variants (SVs) are the predominant class of driver mutation in most cancer types.
Carried only by a subclone, many have a low variant allele frequency, so some studies sequence several samples per patient.
However, current tumor--normal callers need every read of every sample aligned to a single reference genome, so they are computationally expensive and prone to reference bias.
Here we introduce \tool{}, which discovers somatic SVs in a personal pangenome, a colored de Bruijn-like graph of locally consistent parsing (LCP) cores built from a patient's normal and tumor PacBio HiFi reads.
\tool{} first searches this graph, without any alignment, for candidate sites that only tumor reads traverse, and then aligns only short snippets of tumor reads around these sites.
On five CASTLE tumor--normal pairs, \tool{} had the highest mean F1 score against a consensus of 18 call sets (0.881, compared with 0.858 for the next-best caller).
Its lead was largest at the lowest depth and purity (mean F1 score 0.476 versus 0.361 at 5$\times$ and 0.679 versus 0.520 at 20\% purity, each against the next-best caller in that setting).
End to end, \tool{} processed a pair in about 11 minutes on 16 threads, about 24 times faster than mapping the two read sets alone.
For 20 simulated tumors of one patient, a joint run was faster and more accurate than 20 pairwise runs, with the largest accuracy gain at the lowest tumor depth.
Therefore, by calling a patient's samples jointly and losing less accuracy at lower depth and purity, \tool{} can make multi-sample studies more efficient and accurate.
\tool{} is available at \url{https://github.com/Muti-Kara/janus}.
\end{abstract}
